\nivelbasico


\begin{ejercicio}
Un balón se lanza con rapidez inicial:

\[
v_0=\SI{20.0}{\meter\per\second}
\]

formando un ángulo:

\[
\theta_0=30.0\grado
\]

sobre la horizontal.

Determine las componentes rectangulares de la velocidad inicial.

\textit{Pista: descomponga primero el vector velocidad sobre los ejes \(x\) y \(y\).}
\end{ejercicio}

\begin{respuesta}
\[
\boxed{
v_{0x}\approx\SI{17.3}{\meter\per\second}
}
\]

\[
\boxed{
v_{0y}=\SI{10.0}{\meter\per\second}
}
\]
\end{respuesta}

\begin{explicacion}
Tomamos:

\[
+x
\]

horizontal en el sentido del lanzamiento y:

\[
+y
\]

vertical hacia arriba.

El vector velocidad inicial puede escribirse como:

\[
\vect{v}_0
=
v_{0x}\ihat+v_{0y}\jhat.
\]

La componente horizontal es:

\[
v_{0x}
=
v_0\cos\theta_0.
\]

Sustituyendo:

\[
v_{0x}
=
(\SI{20.0}{\meter\per\second})
\cos(30.0\grado).
\]

Entonces:

\[
\boxed{
v_{0x}\approx\SI{17.3}{\meter\per\second}
}.
\]

La componente vertical es:

\[
v_{0y}
=
v_0\sen\theta_0.
\]

Por tanto:

\[
v_{0y}
=
(\SI{20.0}{\meter\per\second})
\sen(30.0\grado).
\]

Así:

\[
\boxed{
v_{0y}=\SI{10.0}{\meter\per\second}
}.
\]

Estas dos componentes describen completamente el vector velocidad inicial:

\[
\vect{v}_0
\approx
(17.3\ihat+10.0\jhat)
\,\si{\meter\per\second}.
\]
\end{explicacion}


\begin{ejercicio}
Una pelota sale horizontalmente del borde de una plataforma situada a:

\[
h=\SI{19.6}{\meter}
\]

sobre el suelo, con rapidez:

\[
v_0=\SI{8.0}{\meter\per\second}.
\]

Despreciando la resistencia del aire, determine:

\begin{enumerate}[label=\alph*)]
    \item el tiempo que tarda en llegar al suelo;
    \item la distancia horizontal recorrida;
    \item las componentes de la velocidad inmediatamente antes del impacto;
    \item la rapidez de impacto.
\end{enumerate}

Use:

\[
g=\SI{9.81}{\meter\per\second\squared}.
\]

\textit{Pista: en un lanzamiento horizontal, \(v_{0y}=0\).}
\end{ejercicio}

\begin{respuesta}
\[
t\approx\SI{2.00}{\second},
\]

\[
\Delta x\approx\SI{16.0}{\meter},
\]

\[
\vect{v}_f
\approx
(8.0\ihat-19.6\jhat)
\,\si{\meter\per\second},
\]

\[
\boxed{
v_f\approx\SI{21.2}{\meter\per\second}
}.
\]
\end{respuesta}

\begin{explicacion}
Elegimos:

\[
+x
\]

horizontal en el sentido del lanzamiento y:

\[
+y
\]

hacia arriba.

Como el lanzamiento es horizontal:

\[
v_{0x}=\SI{8.0}{\meter\per\second},
\]

\[
v_{0y}=0.
\]

Además:

\[
a_x=0,
\qquad
a_y=-g.
\]

Para determinar el tiempo utilizamos el movimiento vertical.

Tomamos el punto de lanzamiento como:

\[
y_0=\SI{19.6}{\meter},
\]

y el suelo como:

\[
y_f=0.
\]

Entonces:

\[
y_f
=
y_0
+
v_{0y}t
-
\frac{1}{2}gt^2.
\]

Como:

\[
v_{0y}=0,
\]

tenemos:

\[
0
=
19.6
-
\frac{1}{2}(9.81)t^2.
\]

Despejando:

\[
t
=
\sqrt{
\frac{2(19.6)}{9.81}
}.
\]

Por tanto:

\[
\boxed{
t\approx\SI{2.00}{\second}
}.
\]

Horizontalmente:

\[
\Delta x=v_{0x}t.
\]

Entonces:

\[
\Delta x
=
(8.0)(2.00).
\]

Así:

\[
\boxed{
\Delta x\approx\SI{16.0}{\meter}
}.
\]

La componente horizontal de la velocidad permanece constante:

\[
\boxed{
v_x=\SI{8.0}{\meter\per\second}
}.
\]

La componente vertical al impactar es:

\[
v_y
=
v_{0y}-gt.
\]

Entonces:

\[
v_y
=
-(9.81)(2.00)
\approx
-\SI{19.6}{\meter\per\second}.
\]

Por tanto:

\[
\boxed{
\vect{v}_f
\approx
(8.0\ihat-19.6\jhat)
\,\si{\meter\per\second}
}.
\]

La rapidez es:

\[
v_f
=
\sqrt{
v_x^2+v_y^2
}.
\]

Entonces:

\[
v_f
=
\sqrt{
(8.0)^2+(19.6)^2
}.
\]

Así:

\[
\boxed{
v_f\approx\SI{21.2}{\meter\per\second}
}.
\]
\end{explicacion}


\begin{ejercicio}
Un proyectil se lanza desde el suelo con:

\[
v_0=\SI{20.0}{\meter\per\second}
\]

y:

\[
\theta_0=30.0\grado.
\]

El proyectil regresa al mismo nivel desde el cual fue lanzado.

Determine:

\begin{enumerate}[label=\alph*)]
    \item el tiempo necesario para alcanzar la altura máxima;
    \item la altura máxima;
    \item el tiempo total de vuelo;
    \item el alcance horizontal.
\end{enumerate}

Use:

\[
g=\SI{9.81}{\meter\per\second\squared}.
\]

\textit{Pista: en la altura máxima, \(v_y=0\).}
\end{ejercicio}

\begin{respuesta}
\[
t_H\approx\SI{1.02}{\second},
\]

\[
H\approx\SI{5.10}{\meter},
\]

\[
t_{\mathrm{vuelo}}\approx\SI{2.04}{\second},
\]

\[
\boxed{
R\approx\SI{35.3}{\meter}
}.
\]
\end{respuesta}

\begin{explicacion}
Primero descomponemos la velocidad inicial:

\[
v_{0x}
=
v_0\cos\theta_0.
\]

Entonces:

\[
v_{0x}
=
20.0\cos(30.0\grado)
\approx
\SI{17.3}{\meter\per\second}.
\]

Verticalmente:

\[
v_{0y}
=
v_0\sen\theta_0.
\]

Por tanto:

\[
v_{0y}
=
20.0\sen(30.0\grado)
=
\SI{10.0}{\meter\per\second}.
\]

En la altura máxima:

\[
v_y=0.
\]

Utilizamos:

\[
v_y=v_{0y}-gt.
\]

Entonces:

\[
0
=
10.0-9.81t_H.
\]

Por tanto:

\[
\boxed{
t_H\approx\SI{1.02}{\second}
}.
\]

La altura máxima respecto del punto de lanzamiento se obtiene mediante:

\[
v_y^2
=
v_{0y}^2-2gH.
\]

Como:

\[
v_y=0,
\]

tenemos:

\[
H
=
\frac{v_{0y}^2}{2g}.
\]

Entonces:

\[
H
=
\frac{(10.0)^2}{2(9.81)}.
\]

Así:

\[
\boxed{
H\approx\SI{5.10}{\meter}
}.
\]

Como el proyectil regresa a la misma altura:

\[
t_{\mathrm{vuelo}}=2t_H.
\]

Por tanto:

\[
\boxed{
t_{\mathrm{vuelo}}
\approx
\SI{2.04}{\second}
}.
\]

Horizontalmente:

\[
R=v_{0x}t_{\mathrm{vuelo}}.
\]

Entonces:

\[
R
=
(17.3)(2.04).
\]

Así:

\[
\boxed{
R\approx\SI{35.3}{\meter}
}.
\]
\end{explicacion}


\begin{ejercicio}
Un proyectil se lanza desde el origen con componentes iniciales:

\[
v_{0x}=\SI{20.0}{\meter\per\second},
\]

\[
v_{0y}=\SI{15.0}{\meter\per\second}.
\]

Determine en:

\[
t=\SI{1.00}{\second}
\]

\begin{enumerate}[label=\alph*)]
    \item su posición;
    \item su vector velocidad;
    \item su rapidez.
\end{enumerate}

Use:

\[
g=\SI{9.81}{\meter\per\second\squared}.
\]
\end{ejercicio}

\begin{respuesta}
\[
\vect{r}
\approx
(20.0\ihat+10.1\jhat)\,\si{\meter},
\]

\[
\vect{v}
=
(20.0\ihat+5.19\jhat)
\,\si{\meter\per\second},
\]

\[
\boxed{
v\approx\SI{20.7}{\meter\per\second}
}.
\]
\end{respuesta}

\begin{explicacion}
Horizontalmente:

\[
a_x=0.
\]

Entonces:

\[
x=v_{0x}t.
\]

Por tanto:

\[
x
=
(20.0)(1.00)
=
\SI{20.0}{\meter}.
\]

Verticalmente:

\[
y
=
v_{0y}t
-
\frac{1}{2}gt^2.
\]

Entonces:

\[
y
=
15.0(1.00)
-
\frac{1}{2}(9.81)(1.00)^2.
\]

Así:

\[
y\approx\SI{10.1}{\meter}.
\]

Por tanto:

\[
\boxed{
\vect{r}
\approx
(20.0\ihat+10.1\jhat)\,\si{\meter}
}.
\]

La componente horizontal de la velocidad permanece constante:

\[
v_x
=
\SI{20.0}{\meter\per\second}.
\]

La componente vertical es:

\[
v_y
=
v_{0y}-gt.
\]

Entonces:

\[
v_y
=
15.0-9.81(1.00).
\]

Así:

\[
v_y
=
\SI{5.19}{\meter\per\second}.
\]

El vector velocidad es:

\[
\boxed{
\vect{v}
=
(20.0\ihat+5.19\jhat)
\,\si{\meter\per\second}
}.
\]

La rapidez es:

\[
v
=
\sqrt{
(20.0)^2+(5.19)^2
}.
\]

Por tanto:

\[
\boxed{
v\approx\SI{20.7}{\meter\per\second}
}.
\]

Como:

\[
v_y>0,
\]

el proyectil todavía está ascendiendo.
\end{explicacion}


\begin{ejercicio}
Un proyectil ideal se encuentra exactamente en el punto más alto de su
trayectoria.

Indique cuáles de las siguientes afirmaciones son correctas:

\begin{enumerate}[label=\alph*)]
    \item \(v_x=0\).
    \item \(v_y=0\).
    \item \(\vect{v}=\vect{0}\).
    \item \(a_x=0\).
    \item \(a_y=-g\).
    \item La aceleración es cero porque el proyectil deja de subir.
\end{enumerate}

Justifique físicamente cada respuesta.
\end{ejercicio}

\begin{respuesta}
Son correctas:

\[
\boxed{\text{b), d) y e)}}.
\]

En general:

\[
v_x\neq0,
\qquad
v_y=0,
\qquad
\vect{a}=-g\jhat.
\]
\end{respuesta}

\begin{explicacion}
En el movimiento ideal de proyectiles:

\[
a_x=0.
\]

Por ello:

\[
v_x=v_{0x}
\]

durante todo el vuelo.

En consecuencia, la afirmación a) es falsa en un lanzamiento oblicuo ordinario.

En el punto más alto:

\[
\boxed{v_y=0}.
\]

Por tanto, b) es correcta.

Sin embargo, como:

\[
v_x\neq0,
\]

el vector velocidad no es nulo:

\[
\vect{v}=v_x\ihat.
\]

Por ello, c) es falsa.

La aceleración horizontal continúa siendo:

\[
\boxed{a_x=0},
\]

de modo que d) es correcta.

La aceleración vertical permanece:

\[
\boxed{a_y=-g}.
\]

Por tanto, e) es correcta.

La afirmación f) es falsa.

El hecho de que:

\[
v_y=0
\]

en un instante no implica:

\[
a_y=0.
\]

De hecho, es precisamente la aceleración gravitatoria la que hace que la
componente vertical cambie de positiva durante el ascenso a negativa durante el
descenso.
\end{explicacion}


\begin{ejercicio}
Desde la misma altura se liberan simultáneamente dos esferas.

La esfera \(A\) se deja caer desde el reposo.

La esfera \(B\) se lanza horizontalmente con rapidez:

\[
v_0=\SI{15}{\meter\per\second}.
\]

Desprecie la resistencia del aire.

\begin{enumerate}[label=\alph*)]
    \item ¿Cuál esfera llega primero al suelo?
    \item ¿Tienen ambas la misma componente vertical de velocidad en cada instante?
    \item ¿Tienen ambas la misma velocidad vectorial?
\end{enumerate}

Explique.
\end{ejercicio}

\begin{respuesta}
\begin{enumerate}[label=\alph*)]
    \item Llegan al suelo simultáneamente.
    \item Sí.
    \item No.
\end{enumerate}
\end{respuesta}

\begin{explicacion}
Ambas esferas comienzan desde la misma altura y tienen:

\[
v_{0y}=0.
\]

También experimentan la misma aceleración vertical:

\[
a_y=-g.
\]

Por tanto, sus movimientos verticales están descritos por exactamente las mismas
ecuaciones:

\[
y
=
y_0-\frac{1}{2}gt^2,
\]

\[
v_y=-gt.
\]

Así, para cualquier instante:

\[
y_A=y_B
\]

y:

\[
(v_y)_A=(v_y)_B.
\]

Por ello:

\[
\boxed{\text{ambas llegan al suelo al mismo tiempo}}.
\]

Sin embargo, la esfera \(B\) posee además una componente horizontal:

\[
(v_x)_B=\SI{15}{\meter\per\second},
\]

mientras que para la esfera \(A\):

\[
(v_x)_A=0.
\]

Por tanto, sus vectores velocidad completos no son iguales.

El movimiento horizontal no altera el movimiento vertical dentro del modelo
ideal.
\end{explicacion}


\begin{ejercicio}
Dos proyectiles se lanzan desde el suelo con la misma rapidez:

\[
v_0=\SI{18.0}{\meter\per\second}.
\]

El primero se lanza a:

\[
\theta_1=30.0\grado,
\]

y el segundo a:

\[
\theta_2=60.0\grado.
\]

Ambos regresan al nivel de lanzamiento.

Determine para cada proyectil:

\begin{enumerate}[label=\alph*)]
    \item el alcance horizontal;
    \item el tiempo de vuelo;
    \item la altura máxima.
\end{enumerate}

Compare los resultados.

\textit{Pista: \(30\grado\) y \(60\grado\) son ángulos complementarios.}
\end{ejercicio}

\begin{respuesta}
Para ambos:

\[
\boxed{
R\approx\SI{28.6}{\meter}
}.
\]

Para \(30\grado\):

\[
t_{\mathrm{vuelo}}\approx\SI{1.83}{\second},
\qquad
H\approx\SI{4.13}{\meter}.
\]

Para \(60\grado\):

\[
t_{\mathrm{vuelo}}\approx\SI{3.18}{\second},
\qquad
H\approx\SI{12.4}{\meter}.
\]
\end{respuesta}

\begin{explicacion}
Para el caso simétrico:

\[
R
=
\frac{
v_0^2\sen(2\theta_0)
}{
g
}.
\]

Para:

\[
\theta_1=30\grado,
\]

tenemos:

\[
\sen(60\grado).
\]

Para:

\[
\theta_2=60\grado,
\]

tenemos:

\[
\sen(120\grado).
\]

Como:

\[
\sen(60\grado)=\sen(120\grado),
\]

ambos alcances son iguales:

\[
R
=
\frac{
(18.0)^2\sen(60\grado)
}{
9.81
}.
\]

Entonces:

\[
\boxed{
R\approx\SI{28.6}{\meter}
}.
\]

Para \(30\grado\):

\[
v_{0y}
=
18.0\sen(30\grado)
=
\SI{9.0}{\meter\per\second}.
\]

El tiempo total es:

\[
t_{\mathrm{vuelo}}
=
\frac{2v_{0y}}{g}.
\]

Entonces:

\[
\boxed{
t_{\mathrm{vuelo}}
\approx
\SI{1.83}{\second}
}.
\]

La altura máxima es:

\[
H
=
\frac{v_{0y}^2}{2g}.
\]

Así:

\[
\boxed{
H\approx\SI{4.13}{\meter}
}.
\]

Para \(60\grado\):

\[
v_{0y}
=
18.0\sen(60\grado)
\approx
\SI{15.6}{\meter\per\second}.
\]

Por tanto:

\[
\boxed{
t_{\mathrm{vuelo}}
\approx
\SI{3.18}{\second}
},
\]

y:

\[
\boxed{
H\approx\SI{12.4}{\meter}
}.
\]

Los dos proyectiles tienen el mismo alcance, pero no siguen la misma trayectoria.

El lanzamiento de \(60\grado\) alcanza mayor altura y permanece más tiempo en el
aire.
\end{explicacion}


\nivelintermedio


\begin{ejercicio}
Desde una plataforma situada a:

\[
h=\SI{15.0}{\meter}
\]

sobre el suelo se lanza una pelota con:

\[
v_0=\SI{18.0}{\meter\per\second}
\]

y:

\[
\theta_0=35.0\grado.
\]

Determine:

\begin{enumerate}[label=\alph*)]
    \item las componentes iniciales de la velocidad;
    \item el tiempo hasta llegar al suelo;
    \item la distancia horizontal recorrida;
    \item el vector velocidad justo antes del impacto;
    \item la rapidez de impacto.
\end{enumerate}

Use:

\[
g=\SI{9.81}{\meter\per\second\squared}.
\]
\end{ejercicio}

\begin{respuesta}
\[
v_{0x}\approx\SI{14.7}{\meter\per\second},
\qquad
v_{0y}\approx\SI{10.3}{\meter\per\second},
\]

\[
t\approx\SI{3.09}{\second},
\]

\[
\Delta x\approx\SI{45.6}{\meter},
\]

\[
\vect{v}_f
\approx
(14.7\ihat-20.0\jhat)
\,\si{\meter\per\second},
\]

\[
\boxed{
v_f\approx\SI{24.9}{\meter\per\second}
}.
\]
\end{respuesta}

\begin{explicacion}
Descomponemos primero la velocidad inicial:

\[
v_{0x}
=
v_0\cos\theta_0.
\]

Entonces:

\[
v_{0x}
=
18.0\cos(35.0\grado)
\approx
\boxed{\SI{14.7}{\meter\per\second}}.
\]

La componente vertical:

\[
v_{0y}
=
18.0\sen(35.0\grado)
\approx
\boxed{\SI{10.3}{\meter\per\second}}.
\]

Tomamos:

\[
y_0=\SI{15.0}{\meter},
\qquad
y_f=0.
\]

La ecuación vertical es:

\[
y_f
=
y_0+v_{0y}t-\frac{1}{2}gt^2.
\]

Por tanto:

\[
0
=
15.0+10.3t-4.905t^2.
\]

Reordenando:

\[
4.905t^2-10.3t-15.0=0.
\]

La ecuación produce una raíz positiva y una negativa.

La raíz físicamente válida es:

\[
\boxed{
t\approx\SI{3.09}{\second}
}.
\]

Horizontalmente:

\[
\Delta x=v_{0x}t.
\]

Entonces:

\[
\Delta x
=
(14.7)(3.09).
\]

Así:

\[
\boxed{
\Delta x\approx\SI{45.6}{\meter}
}.
\]

La componente horizontal de la velocidad al impactar es:

\[
v_x=v_{0x}
\approx
\SI{14.7}{\meter\per\second}.
\]

La componente vertical:

\[
v_y
=
v_{0y}-gt.
\]

Entonces:

\[
v_y
=
10.3-9.81(3.09).
\]

Por tanto:

\[
\boxed{
v_y\approx-\SI{20.0}{\meter\per\second}
}.
\]

Así:

\[
\boxed{
\vect{v}_f
\approx
(14.7\ihat-20.0\jhat)
\,\si{\meter\per\second}
}.
\]

La rapidez de impacto es:

\[
v_f
=
\sqrt{
(14.7)^2+(20.0)^2
}.
\]

Entonces:

\[
\boxed{
v_f\approx\SI{24.9}{\meter\per\second}
}.
\]

La rapidez final es mayor que la inicial porque el proyectil termina a una altura
menor que aquella desde la que fue lanzado.
\end{explicacion}


\begin{ejercicio}
Un helicóptero vuela horizontalmente a:

\[
\SI{72.0}{\kilo\meter\per\hour}
\]

a una altura de:

\[
\SI{180}{\meter}.
\]

Debe liberar un paquete para que llegue a un punto situado en el suelo.

Despreciando la resistencia del aire, determine:

\begin{enumerate}[label=\alph*)]
    \item la rapidez horizontal en unidades SI;
    \item el tiempo de caída;
    \item a qué distancia horizontal antes del punto debe liberarse el paquete;
    \item las componentes de la velocidad al impactar;
    \item la rapidez de impacto.
\end{enumerate}
\end{ejercicio}

\begin{respuesta}
\[
v_x=\SI{20.0}{\meter\per\second},
\]

\[
t\approx\SI{6.06}{\second},
\]

\[
\Delta x\approx\SI{121}{\meter},
\]

\[
\vect{v}_f
\approx
(20.0\ihat-59.4\jhat)
\,\si{\meter\per\second},
\]

\[
\boxed{
v_f\approx\SI{62.7}{\meter\per\second}
}.
\]
\end{respuesta}

\begin{explicacion}
Primero convertimos la rapidez:

\[
\SI{72.0}{\kilo\meter\per\hour}
\left(
\frac{\SI{1000}{\meter}}{\SI{1}{\kilo\meter}}
\right)
\left(
\frac{\SI{1}{\hour}}{\SI{3600}{\second}}
\right).
\]

Así:

\[
\boxed{
v_x=\SI{20.0}{\meter\per\second}
}.
\]

Al liberarse, el paquete conserva la velocidad horizontal del helicóptero.

Verticalmente:

\[
v_{0y}=0.
\]

La caída desde:

\[
h=\SI{180}{\meter}
\]

satisface:

\[
h=\frac{1}{2}gt^2.
\]

Entonces:

\[
t
=
\sqrt{
\frac{2h}{g}
}.
\]

Sustituyendo:

\[
t
=
\sqrt{
\frac{2(180)}{9.81}
}.
\]

Así:

\[
\boxed{
t\approx\SI{6.06}{\second}
}.
\]

Horizontalmente:

\[
\Delta x=v_xt.
\]

Entonces:

\[
\Delta x
=
(20.0)(6.06).
\]

Por tanto:

\[
\boxed{
\Delta x\approx\SI{121}{\meter}
}.
\]

La componente vertical de la velocidad al impactar es:

\[
v_y=-gt.
\]

Así:

\[
v_y
=
-9.81(6.06)
\approx
-\SI{59.4}{\meter\per\second}.
\]

El vector velocidad es:

\[
\boxed{
\vect{v}_f
\approx
(20.0\ihat-59.4\jhat)
\,\si{\meter\per\second}
}.
\]

La rapidez:

\[
v_f
=
\sqrt{
20.0^2+59.4^2
}.
\]

Entonces:

\[
\boxed{
v_f\approx\SI{62.7}{\meter\per\second}
}.
\]
\end{explicacion}


\begin{ejercicio}
Un proyectil parte de:

\[
(x_0,y_0)=(0,\SI{2.0}{\meter})
\]

con componentes iniciales:

\[
v_{0x}=\SI{15.0}{\meter\per\second},
\]

\[
v_{0y}=\SI{20.0}{\meter\per\second}.
\]

Determine:

\begin{enumerate}[label=\alph*)]
    \item la ecuación de la trayectoria \(y(x)\);
    \item el tiempo para alcanzar la altura máxima;
    \item las coordenadas del punto más alto;
    \item la posición horizontal donde llega al suelo.
\end{enumerate}
\end{ejercicio}

\begin{respuesta}
\[
\boxed{
y
=
2.0
+
1.333x
-
0.0218x^2
}
\]

con \(x\) y \(y\) en metros,

\[
t_H\approx\SI{2.04}{\second},
\]

\[
(x_H,y_H)
\approx
(\SI{30.6}{\meter},\SI{22.4}{\meter}),
\]

\[
\boxed{
x_f\approx\SI{62.6}{\meter}
}.
\]
\end{respuesta}

\begin{explicacion}
Las ecuaciones paramétricas son:

\[
x=v_{0x}t,
\]

\[
y
=
y_0+v_{0y}t-\frac{1}{2}gt^2.
\]

De la primera:

\[
t=\frac{x}{v_{0x}}.
\]

Entonces:

\[
t=\frac{x}{15.0}.
\]

Sustituimos en la ecuación vertical:

\[
y
=
2.0
+
20.0
\left(
\frac{x}{15.0}
\right)
-
\frac{1}{2}(9.81)
\left(
\frac{x}{15.0}
\right)^2.
\]

Simplificando:

\[
\boxed{
y
=
2.0
+
1.333x
-
0.0218x^2
}.
\]

El término cuadrático negativo confirma que la parábola abre hacia abajo.

Para la altura máxima:

\[
v_y=0.
\]

Entonces:

\[
0=20.0-9.81t_H.
\]

Así:

\[
\boxed{
t_H\approx\SI{2.04}{\second}
}.
\]

La coordenada horizontal correspondiente es:

\[
x_H
=
v_{0x}t_H.
\]

Entonces:

\[
x_H
=
15.0(2.04)
\approx
\boxed{\SI{30.6}{\meter}}.
\]

La altura máxima es:

\[
y_H
=
y_0+\frac{v_{0y}^2}{2g}.
\]

Entonces:

\[
y_H
=
2.0
+
\frac{(20.0)^2}{2(9.81)}.
\]

Por tanto:

\[
\boxed{
y_H\approx\SI{22.4}{\meter}
}.
\]

Para encontrar dónde llega al suelo hacemos:

\[
y=0.
\]

Entonces:

\[
0
=
2.0
+
1.333x
-
0.0218x^2.
\]

La ecuación cuadrática proporciona dos soluciones.

Una es negativa y corresponde a la prolongación matemática de la parábola antes
del lanzamiento.

La solución física posterior al lanzamiento es:

\[
\boxed{
x_f\approx\SI{62.6}{\meter}
}.
\]
\end{explicacion}


\begin{ejercicio}
Un proyectil se lanza desde el origen con:

\[
v_0=\SI{24.0}{\meter\per\second}
\]

y:

\[
\theta_0=40.0\grado.
\]

Determine los dos instantes en que el proyectil se encuentra a:

\[
y=\SI{8.0}{\meter}.
\]

Para cada instante, determine:

\begin{enumerate}[label=\alph*)]
    \item la posición horizontal;
    \item la componente vertical de la velocidad;
    \item si el proyectil asciende o desciende.
\end{enumerate}
\end{ejercicio}

\begin{respuesta}
Primer paso por la altura:

\[
t_1\approx\SI{0.655}{\second},
\]

\[
x_1\approx\SI{12.0}{\meter},
\]

\[
v_{y1}\approx+\SI{9.00}{\meter\per\second}.
\]

Segundo paso:

\[
t_2\approx\SI{2.49}{\second},
\]

\[
x_2\approx\SI{45.8}{\meter},
\]

\[
v_{y2}\approx-\SI{9.00}{\meter\per\second}.
\]
\end{respuesta}

\begin{explicacion}
Las componentes iniciales son:

\[
v_{0x}
=
24.0\cos(40.0\grado)
\approx
\SI{18.4}{\meter\per\second},
\]

\[
v_{0y}
=
24.0\sen(40.0\grado)
\approx
\SI{15.4}{\meter\per\second}.
\]

La ecuación vertical es:

\[
y
=
v_{0y}t-\frac{1}{2}gt^2.
\]

Para:

\[
y=\SI{8.0}{\meter},
\]

tenemos:

\[
8.0
=
15.4t-4.905t^2.
\]

Reordenando:

\[
4.905t^2-15.4t+8.0=0.
\]

La ecuación cuadrática produce:

\[
\boxed{
t_1\approx\SI{0.655}{\second}
}
\]

y:

\[
\boxed{
t_2\approx\SI{2.49}{\second}
}.
\]

Las posiciones horizontales son:

\[
x_1=v_{0x}t_1.
\]

Entonces:

\[
\boxed{
x_1\approx\SI{12.0}{\meter}
}.
\]

Para el segundo instante:

\[
x_2=v_{0x}t_2.
\]

Así:

\[
\boxed{
x_2\approx\SI{45.8}{\meter}
}.
\]

La componente vertical de velocidad es:

\[
v_y=v_{0y}-gt.
\]

Para el primer instante:

\[
v_{y1}
=
15.4-9.81(0.655)
\approx
+\SI{9.00}{\meter\per\second}.
\]

Por tanto, el proyectil está ascendiendo.

Para el segundo:

\[
v_{y2}
=
15.4-9.81(2.49)
\approx
-\SI{9.00}{\meter\per\second}.
\]

Por tanto, está descendiendo.

Las magnitudes son iguales y los signos opuestos porque se trata de la misma
altura en una trayectoria ideal simétrica respecto del ápice.
\end{explicacion}


\begin{ejercicio}
Una pelota se lanza desde el suelo con:

\[
v_0=\SI{22.0}{\meter\per\second}
\]

y:

\[
\theta_0=45.0\grado.
\]

A:

\[
\SI{20.0}{\meter}
\]

del punto de lanzamiento existe una pared vertical de:

\[
\SI{10.0}{\meter}
\]

de altura.

Determine:

\begin{enumerate}[label=\alph*)]
    \item la altura de la pelota cuando llega a la posición de la pared;
    \item si supera la pared;
    \item las componentes de la velocidad en ese instante;
    \item la rapidez.
\end{enumerate}
\end{ejercicio}

\begin{respuesta}
\[
y\approx\SI{11.9}{\meter}.
\]

La pelota supera la pared por aproximadamente:

\[
\SI{1.89}{\meter}.
\]

\[
\vect{v}
\approx
(15.6\ihat+2.94\jhat)
\,\si{\meter\per\second},
\]

\[
\boxed{
v\approx\SI{15.8}{\meter\per\second}
}.
\]
\end{respuesta}

\begin{explicacion}
Las componentes iniciales son:

\[
v_{0x}
=
22.0\cos(45.0\grado)
\approx
\SI{15.6}{\meter\per\second},
\]

\[
v_{0y}
=
22.0\sen(45.0\grado)
\approx
\SI{15.6}{\meter\per\second}.
\]

La pelota llega a la pared cuando:

\[
x=\SI{20.0}{\meter}.
\]

Horizontalmente:

\[
x=v_{0x}t.
\]

Entonces:

\[
t
=
\frac{20.0}{15.6}
\approx
\SI{1.29}{\second}.
\]

La altura es:

\[
y
=
v_{0y}t-\frac{1}{2}gt^2.
\]

Sustituyendo:

\[
y
=
15.6(1.29)
-
4.905(1.29)^2.
\]

Por tanto:

\[
\boxed{
y\approx\SI{11.9}{\meter}
}.
\]

Como la pared mide:

\[
\SI{10.0}{\meter},
\]

la diferencia es:

\[
11.9-10.0
=
\SI{1.9}{\meter}.
\]

Por tanto, la pelota sí supera la pared.

La velocidad horizontal sigue siendo:

\[
v_x\approx\SI{15.6}{\meter\per\second}.
\]

Verticalmente:

\[
v_y
=
v_{0y}-gt.
\]

Así:

\[
v_y
=
15.6-9.81(1.29)
\approx
+\SI{2.94}{\meter\per\second}.
\]

El signo positivo indica que todavía está ascendiendo.

Entonces:

\[
\boxed{
\vect{v}
\approx
(15.6\ihat+2.94\jhat)
\,\si{\meter\per\second}
}.
\]

La rapidez es:

\[
v
=
\sqrt{
15.6^2+2.94^2
}.
\]

Así:

\[
\boxed{
v\approx\SI{15.8}{\meter\per\second}
}.
\]
\end{explicacion}


\begin{ejercicio}
Un proyectil se lanza desde el origen. Después de:

\[
t=\SI{2.00}{\second}
\]

se encuentra en:

\[
\vect{r}
=
(30.0\ihat+10.0\jhat)\,\si{\meter}.
\]

Despreciando la resistencia del aire, determine:

\begin{enumerate}[label=\alph*)]
    \item las componentes iniciales de la velocidad;
    \item la rapidez inicial;
    \item el ángulo inicial de lanzamiento;
    \item el vector velocidad en \(t=\SI{2.00}{\second}\);
    \item si el proyectil está ascendiendo o descendiendo en ese instante.
\end{enumerate}
\end{ejercicio}

\begin{respuesta}
\[
v_{0x}=\SI{15.0}{\meter\per\second},
\]

\[
v_{0y}\approx\SI{14.8}{\meter\per\second},
\]

\[
v_0\approx\SI{21.1}{\meter\per\second},
\]

\[
\theta_0\approx44.7\grado,
\]

\[
\vect{v}(2.00)
\approx
(15.0\ihat-4.81\jhat)
\,\si{\meter\per\second}.
\]

El proyectil está descendiendo.
\end{respuesta}

\begin{explicacion}
Horizontalmente:

\[
x=v_{0x}t.
\]

Por tanto:

\[
v_{0x}
=
\frac{x}{t}.
\]

Sustituyendo:

\[
v_{0x}
=
\frac{\SI{30.0}{\meter}}
{\SI{2.00}{\second}}.
\]

Así:

\[
\boxed{
v_{0x}=\SI{15.0}{\meter\per\second}
}.
\]

Verticalmente:

\[
y
=
v_{0y}t-\frac{1}{2}gt^2.
\]

Entonces:

\[
10.0
=
v_{0y}(2.00)
-
\frac{1}{2}(9.81)(2.00)^2.
\]

Por tanto:

\[
10.0
=
2v_{0y}-19.62.
\]

Despejando:

\[
\boxed{
v_{0y}\approx\SI{14.8}{\meter\per\second}
}.
\]

La rapidez inicial es:

\[
v_0
=
\sqrt{
v_{0x}^2+v_{0y}^2
}.
\]

Entonces:

\[
v_0
=
\sqrt{
(15.0)^2+(14.8)^2
}.
\]

Así:

\[
\boxed{
v_0\approx\SI{21.1}{\meter\per\second}
}.
\]

La dirección inicial satisface:

\[
\tan\theta_0
=
\frac{v_{0y}}{v_{0x}}.
\]

Por tanto:

\[
\boxed{
\theta_0\approx44.7\grado
}.
\]

En \(t=\SI{2.00}{\second}\):

\[
v_x=v_{0x}
=
\SI{15.0}{\meter\per\second}.
\]

Verticalmente:

\[
v_y
=
v_{0y}-gt.
\]

Entonces:

\[
v_y
=
14.8-9.81(2.00)
\approx
-\SI{4.81}{\meter\per\second}.
\]

Por tanto:

\[
\boxed{
\vect{v}(2.00)
\approx
(15.0\ihat-4.81\jhat)
\,\si{\meter\per\second}
}.
\]

Como:

\[
v_y<0,
\]

el proyectil está descendiendo.
\end{explicacion}


\nivelavanzado


\begin{ejercicio}
Un lanzador puede impartir a una pelota una rapidez inicial fija de:

\[
v_0=\SI{30.0}{\meter\per\second}.
\]

Se desea que la pelota pase por el punto:

\[
(x,y)
=
(\SI{50.0}{\meter},\SI{10.0}{\meter})
\]

respecto del punto de lanzamiento.

Determine:

\begin{enumerate}[label=\alph*)]
    \item todos los ángulos de lanzamiento posibles entre \(0\grado\) y
    \(90\grado\);
    \item el tiempo necesario para alcanzar el punto en cada trayectoria;
    \item cuál trayectoria es más alta.
\end{enumerate}

Desprecie la resistencia del aire.
\end{ejercicio}

\begin{respuesta}
Existen dos ángulos:

\[
\boxed{
\theta_1\approx29.1\grado
}
\]

y:

\[
\boxed{
\theta_2\approx72.2\grado
}.
\]

Los tiempos son aproximadamente:

\[
t_1\approx\SI{1.91}{\second},
\]

\[
t_2\approx\SI{5.45}{\second}.
\]

La trayectoria asociada a \(72.2\grado\) es la más alta.
\end{respuesta}

\begin{explicacion}
La ecuación general de la trayectoria, tomando el lanzamiento desde el origen,
es:

\[
y
=
x\tan\theta_0
-
\frac{
gx^2
}{
2v_0^2\cos^2\theta_0
}.
\]

Introducimos:

\[
T=\tan\theta_0.
\]

Utilizamos:

\[
\frac{1}{\cos^2\theta_0}
=
1+\tan^2\theta_0
=
1+T^2.
\]

Entonces:

\[
y
=
xT
-
\frac{gx^2}{2v_0^2}
(1+T^2).
\]

Sustituimos:

\[
x=\SI{50.0}{\meter},
\]

\[
y=\SI{10.0}{\meter},
\]

\[
v_0=\SI{30.0}{\meter\per\second}.
\]

Calculamos:

\[
\frac{gx^2}{2v_0^2}
=
\frac{
9.81(50.0)^2
}{
2(30.0)^2
}
=
13.625.
\]

Por tanto:

\[
10
=
50T
-
13.625(1+T^2).
\]

Reordenando:

\[
13.625T^2
-
50T
+
23.625
=
0.
\]

Aplicando la fórmula cuadrática:

\[
T
=
\frac{
50
\pm
\sqrt{
50^2
-
4(13.625)(23.625)
}
}{
2(13.625)
}.
\]

Se obtienen:

\[
T_1\approx0.557,
\]

\[
T_2\approx3.113.
\]

Entonces:

\[
\theta_1
=
\arctan(0.557)
\approx
\boxed{29.1\grado},
\]

\[
\theta_2
=
\arctan(3.113)
\approx
\boxed{72.2\grado}.
\]

Para cada trayectoria:

\[
x
=
v_0\cos\theta_0\,t.
\]

Por tanto:

\[
t
=
\frac{x}
{v_0\cos\theta_0}.
\]

Para el ángulo menor:

\[
\boxed{
t_1\approx\SI{1.91}{\second}
}.
\]

Para el ángulo mayor:

\[
\boxed{
t_2\approx\SI{5.45}{\second}
}.
\]

Existen dos trayectorias posibles para alcanzar el mismo punto: una trayectoria
baja y rápida y otra más alta con mayor tiempo de vuelo.

La correspondiente a:

\[
72.2\grado
\]

posee una componente vertical inicial mucho mayor y, por tanto, alcanza una
altura máxima mayor.
\end{explicacion}


\begin{ejercicio}
Para un proyectil lanzado desde el suelo con rapidez inicial \(v_0\) y ángulo
\(\theta_0\), que regresa al mismo nivel, deduzca:

\begin{enumerate}[label=\alph*)]
    \item el tiempo total de vuelo;
    \item el alcance horizontal;
    \item la relación entre los alcances de dos ángulos complementarios;
    \item el ángulo que produce el alcance máximo.
\end{enumerate}

Indique claramente las condiciones bajo las cuales los resultados son válidos.
\end{ejercicio}

\begin{respuesta}
\[
\boxed{
t_{\mathrm{vuelo}}
=
\frac{
2v_0\sen\theta_0
}{
g
}
},
\]

\[
\boxed{
R
=
\frac{
v_0^2\sen(2\theta_0)
}{
g
}
}.
\]

Los ángulos complementarios producen el mismo alcance.

El alcance máximo ocurre en:

\[
\boxed{
\theta_0=45\grado
}.
\]
\end{respuesta}

\begin{explicacion}
Las componentes iniciales son:

\[
v_{0x}
=
v_0\cos\theta_0,
\]

\[
v_{0y}
=
v_0\sen\theta_0.
\]

Como el proyectil regresa al nivel inicial:

\[
y_f-y_0=0.
\]

La ecuación vertical es:

\[
0
=
v_{0y}t
-
\frac{1}{2}gt^2.
\]

Factorizamos:

\[
0
=
t
\left(
v_{0y}
-
\frac{1}{2}gt
\right).
\]

Una solución es:

\[
t=0,
\]

que corresponde al instante de lanzamiento.

La segunda solución corresponde al regreso:

\[
v_{0y}
-
\frac{1}{2}gt
=
0.
\]

Por tanto:

\[
t_{\mathrm{vuelo}}
=
\frac{
2v_{0y}
}{
g
}.
\]

Sustituyendo:

\[
v_{0y}
=
v_0\sen\theta_0,
\]

obtenemos:

\[
\boxed{
t_{\mathrm{vuelo}}
=
\frac{
2v_0\sen\theta_0
}{
g
}
}.
\]

Horizontalmente:

\[
R
=
v_{0x}t_{\mathrm{vuelo}}.
\]

Entonces:

\[
R
=
v_0\cos\theta_0
\left(
\frac{
2v_0\sen\theta_0
}{
g
}
\right).
\]

Por tanto:

\[
R
=
\frac{
2v_0^2
\sen\theta_0
\cos\theta_0
}{
g
}.
\]

Utilizando:

\[
2\sen\theta_0\cos\theta_0
=
\sen(2\theta_0),
\]

resulta:

\[
\boxed{
R
=
\frac{
v_0^2\sen(2\theta_0)
}{
g
}
}.
\]

Ahora consideremos:

\[
\theta_2
=
90\grado-\theta_1.
\]

Entonces:

\[
2\theta_2
=
180\grado-2\theta_1.
\]

Como:

\[
\sen(180\grado-\alpha)
=
\sen\alpha,
\]

tenemos:

\[
R(\theta_1)
=
R(\theta_2).
\]

Por tanto, los ángulos complementarios producen el mismo alcance.

Para rapidez inicial fija:

\[
R
\propto
\sen(2\theta_0).
\]

El valor máximo del seno es:

\[
1.
\]

Entonces:

\[
2\theta_0=90\grado.
\]

Por tanto:

\[
\boxed{
\theta_0=45\grado
}.
\]

Estos resultados requieren:

\begin{itemize}
    \item resistencia del aire despreciable;
    \item \(g\) constante;
    \item misma altura inicial y final;
    \item superficie de referencia horizontal.
\end{itemize}

Fuera de estas condiciones, el resultado de \(45\grado\) no es necesariamente
válido.
\end{explicacion}


\begin{ejercicio}
Un proyectil se lanza desde una plataforma situada a una altura \(h\) sobre un
terreno horizontal con rapidez inicial fija \(v_0\).

Despreciando la resistencia del aire:

\begin{enumerate}[label=\alph*)]
    \item deduzca una expresión para el ángulo de lanzamiento que maximiza el
    alcance horizontal;
    \item deduzca una expresión para el alcance máximo;
    \item evalúe ambos resultados para:
    \[
    v_0=\SI{25.0}{\meter\per\second},
    \qquad
    h=\SI{20.0}{\meter}.
    \]
\end{enumerate}

Use:

\[
g=\SI{9.81}{\meter\per\second\squared}.
\]
\end{ejercicio}

\begin{respuesta}
\[
\boxed{
\tan\theta_{\max}
=
\frac{
v_0
}{
\sqrt{v_0^2+2gh}
}
}
\]

y:

\[
\boxed{
R_{\max}
=
\frac{v_0}{g}
\sqrt{
v_0^2+2gh
}
}.
\]

Para los datos:

\[
\boxed{
\theta_{\max}\approx38.1\grado
}
\]

y:

\[
\boxed{
R_{\max}\approx\SI{81.3}{\meter}
}.
\]
\end{respuesta}

\begin{explicacion}
Tomamos el punto de lanzamiento como:

\[
(x_0,y_0)=(0,h).
\]

La ecuación de la trayectoria es:

\[
y
=
h
+
x\tan\theta
-
\frac{
gx^2
}{
2v_0^2\cos^2\theta
}.
\]

Al llegar al suelo:

\[
y=0,
\qquad
x=R.
\]

Entonces:

\[
0
=
h
+
R\tan\theta
-
\frac{
gR^2
}{
2v_0^2\cos^2\theta
}.
\]

Definimos:

\[
T=\tan\theta.
\]

Como:

\[
\frac{1}{\cos^2\theta}
=
1+T^2,
\]

obtenemos:

\[
0
=
h
+
RT
-
\frac{
gR^2
}{
2v_0^2
}
(1+T^2).
\]

Definamos:

\[
K
=
\frac{
gR^2
}{
2v_0^2
}.
\]

Entonces:

\[
KT^2
-
RT
+
(K-h)
=
0.
\]

Para un alcance \(R\) dado, esta ecuación debe tener una solución real para
\(T\).

El alcance máximo corresponde al caso límite en el que las dos posibles
soluciones para el ángulo coinciden.

Por tanto, el discriminante debe ser cero:

\[
R^2
-
4K(K-h)
=
0.
\]

Sustituimos:

\[
K
=
\frac{
gR^2
}{
2v_0^2
}.
\]

Después de simplificar:

\[
R_{\max}^2
=
\frac{
v_0^2
}{
g^2
}
\left(
v_0^2+2gh
\right).
\]

Por tanto:

\[
\boxed{
R_{\max}
=
\frac{v_0}{g}
\sqrt{
v_0^2+2gh
}
}.
\]

Cuando el discriminante es cero, la raíz de la ecuación cuadrática es:

\[
T
=
\frac{R}{2K}.
\]

Sustituyendo el valor de \(K\):

\[
T
=
\frac{
v_0^2
}{
gR_{\max}
}.
\]

Utilizando la expresión obtenida para \(R_{\max}\):

\[
\boxed{
\tan\theta_{\max}
=
\frac{
v_0
}{
\sqrt{
v_0^2+2gh
}
}
}.
\]

Ahora sustituimos:

\[
v_0=\SI{25.0}{\meter\per\second},
\]

\[
h=\SI{20.0}{\meter}.
\]

Entonces:

\[
\sqrt{
v_0^2+2gh
}
=
\sqrt{
25.0^2+2(9.81)(20.0)
}.
\]

Esto da aproximadamente:

\[
\SI{31.9}{\meter\per\second}.
\]

Así:

\[
\tan\theta_{\max}
=
\frac{25.0}{31.9}.
\]

Por tanto:

\[
\boxed{
\theta_{\max}\approx38.1\grado
}.
\]

El alcance máximo es:

\[
R_{\max}
=
\frac{25.0}{9.81}(31.9).
\]

Entonces:

\[
\boxed{
R_{\max}\approx\SI{81.3}{\meter}
}.
\]

Observe que:

\[
38.1\grado<45\grado.
\]

Esto confirma que \(45\grado\) no es universalmente el ángulo de alcance máximo.
Cuando el punto de lanzamiento está por encima del punto de llegada, el
proyectil dispone de tiempo adicional de caída y el ángulo óptimo es menor.
\end{explicacion}


\begin{ejercicio}
En un experimento se lanza horizontalmente una pequeña esfera.

El análisis de video muestra que su trayectoria, tomando como origen el punto de
lanzamiento y \(+y\) hacia arriba, está bien descrita por:

\[
y=-0.0300x^2,
\]

donde \(x\) y \(y\) se miden en metros.

La componente horizontal de la velocidad, medida experimentalmente, es:

\[
v_x=\SI{12.8}{\meter\per\second}.
\]

Suponiendo que la resistencia del aire es despreciable:

\begin{enumerate}[label=\alph*)]
    \item deduzca la relación entre el coeficiente de \(x^2\), \(g\) y \(v_x\);
    \item estime experimentalmente \(g\);
    \item compare el resultado con:
    \[
    g=\SI{9.81}{\meter\per\second\squared}.
    \]
\end{enumerate}
\end{ejercicio}

\begin{respuesta}
\[
y
=
-\frac{g}{2v_x^2}x^2,
\]

\[
\boxed{
g_{\mathrm{exp}}
\approx
\SI{9.83}{\meter\per\second\squared}
}.
\]

El resultado es muy cercano al valor de referencia.
\end{respuesta}

\begin{explicacion}
Para un lanzamiento horizontal desde el origen:

\[
x=v_xt.
\]

Por tanto:

\[
t=\frac{x}{v_x}.
\]

Verticalmente:

\[
y=-\frac{1}{2}gt^2.
\]

Sustituimos el tiempo:

\[
y
=
-\frac{1}{2}g
\left(
\frac{x}{v_x}
\right)^2.
\]

Entonces:

\[
\boxed{
y
=
-\frac{g}{2v_x^2}x^2
}.
\]

Comparando con el ajuste experimental:

\[
y=-0.0300x^2,
\]

obtenemos:

\[
0.0300
=
\frac{g}{2v_x^2}.
\]

Despejando:

\[
g
=
2(0.0300)v_x^2.
\]

Sustituyendo:

\[
v_x=\SI{12.8}{\meter\per\second},
\]

tenemos:

\[
g
=
2(0.0300)(12.8)^2.
\]

Entonces:

\[
\boxed{
g_{\mathrm{exp}}
\approx
\SI{9.83}{\meter\per\second\squared}
}.
\]

Comparado con:

\[
\SI{9.81}{\meter\per\second\squared},
\]

la diferencia es pequeña.

El experimento muestra que la forma de la trayectoria puede utilizarse para
inferir cuantitativamente la aceleración gravitatoria cuando el modelo ideal es
una buena aproximación.
\end{explicacion}


\begin{ejercicio}
La trayectoria experimental idealizada de un proyectil está descrita por:

\[
y
=
2.00
+
0.750x
-
0.0125x^2,
\]

donde \(x\) y \(y\) se expresan en metros.

Suponiendo movimiento ideal de proyectil con:

\[
g=\SI{9.81}{\meter\per\second\squared},
\]

determine:

\begin{enumerate}[label=\alph*)]
    \item la altura inicial;
    \item el ángulo inicial de lanzamiento;
    \item la rapidez inicial;
    \item las componentes iniciales de la velocidad;
    \item las coordenadas de la altura máxima.
\end{enumerate}
\end{ejercicio}

\begin{respuesta}
\[
y_0=\SI{2.00}{\meter},
\]

\[
\theta_0\approx36.9\grado,
\]

\[
v_0\approx\SI{24.8}{\meter\per\second},
\]

\[
v_{0x}\approx\SI{19.8}{\meter\per\second},
\]

\[
v_{0y}\approx\SI{14.9}{\meter\per\second},
\]

\[
\boxed{
(x_H,y_H)
=
(\SI{30.0}{\meter},\SI{13.25}{\meter})
}.
\]
\end{respuesta}

\begin{explicacion}
La ecuación general de la trayectoria es:

\[
y
=
y_0
+
x\tan\theta_0
-
\frac{
gx^2
}{
2v_0^2\cos^2\theta_0
}.
\]

La trayectoria dada es:

\[
y
=
2.00
+
0.750x
-
0.0125x^2.
\]

Comparando términos constantes:

\[
\boxed{
y_0=\SI{2.00}{\meter}
}.
\]

Comparando los términos lineales:

\[
\tan\theta_0=0.750.
\]

Por tanto:

\[
\theta_0
=
\arctan(0.750).
\]

Así:

\[
\boxed{
\theta_0\approx36.9\grado
}.
\]

Para este ángulo:

\[
\cos\theta_0=0.800,
\]

por lo que:

\[
\cos^2\theta_0=0.640.
\]

Comparamos ahora el coeficiente cuadrático:

\[
0.0125
=
\frac{
g
}{
2v_0^2\cos^2\theta_0
}.
\]

Entonces:

\[
v_0^2
=
\frac{
g
}{
2(0.0125)(0.640)
}.
\]

Sustituyendo:

\[
v_0^2
=
\frac{
9.81
}{
0.0160
}.
\]

Por tanto:

\[
\boxed{
v_0\approx\SI{24.8}{\meter\per\second}
}.
\]

Las componentes son:

\[
v_{0x}
=
v_0\cos\theta_0.
\]

Así:

\[
\boxed{
v_{0x}\approx\SI{19.8}{\meter\per\second}
}.
\]

La componente vertical:

\[
v_{0y}
=
v_0\sen\theta_0.
\]

Entonces:

\[
\boxed{
v_{0y}\approx\SI{14.9}{\meter\per\second}
}.
\]

Para hallar el punto máximo podemos utilizar la propia parábola:

\[
y
=
2.00+0.750x-0.0125x^2.
\]

El vértice de:

\[
y=ax^2+bx+c
\]

se encuentra en:

\[
x_H
=
-\frac{b}{2a}.
\]

En este caso:

\[
a=-0.0125,
\qquad
b=0.750.
\]

Entonces:

\[
x_H
=
-\frac{0.750}{2(-0.0125)}
=
\boxed{\SI{30.0}{\meter}}.
\]

Sustituyendo en la trayectoria:

\[
y_H
=
2.00
+
0.750(30.0)
-
0.0125(30.0)^2.
\]

Entonces:

\[
y_H
=
2.00+22.5-11.25.
\]

Por tanto:

\[
\boxed{
y_H=\SI{13.25}{\meter}
}.
\]

Así:

\[
\boxed{
(x_H,y_H)
=
(\SI{30.0}{\meter},\SI{13.25}{\meter})
}.
\]

Este problema muestra que la trayectoria espacial contiene información suficiente
para reconstruir varias características del lanzamiento.
\end{explicacion}